\documentclass[11pt,letterpaper]{article}
\newcommand{\labid}{1-E}
\newcommand{\labtitleshort}{Vivado Loopback}
% Shared preamble for Module 1 lab handouts.
% Usage: each handout defines \labid and \labtitleshort, then % Shared preamble for Module 1 lab handouts.
% Usage: each handout defines \labid and \labtitleshort, then % Shared preamble for Module 1 lab handouts.
% Usage: each handout defines \labid and \labtitleshort, then \input{LabHandout_Preamble}
% BEFORE \begin{document}. Compiles with a single pdflatex pass (no glossaries tool).

\usepackage[T1]{fontenc}
\usepackage[margin=1in]{geometry}
\usepackage{titlesec}
\usepackage{enumitem}
\usepackage{booktabs}
\usepackage{tabularx}
\usepackage{graphicx}
\graphicspath{{Images/}{../Images/}}
\usepackage{xcolor}
\usepackage{hyperref}
\usepackage{fancyhdr}
\usepackage{amsmath}
\usepackage{amssymb}
\usepackage{tcolorbox}
\usepackage{caption}
\tcbuselibrary{skins,breakable}

% ── Glossary (per-document, built from shared glossary.tex) ──
% Uses \makenoidxglossaries so no external makeglossaries tool is
% needed: sorting and printing is done by LaTeX itself across two
% pdflatex passes (same cadence as the TOC). Only acronyms that
% are actually \gls{}-used in the handout appear in the printed
% glossary, and each \gls{} reference is a live hyperlink that
% jumps to its entry in the Acronyms section at the end.
\usepackage[acronym,nomain,nopostdot]{glossaries}
\makenoidxglossaries
\input{../../glossary}
% Call \printhandoutglossary just before \end{document} in each handout.
\newcommand{\printhandoutglossary}{%
	\clearpage
	\printnoidxglossary[type=\acronymtype,title={Glossary}]%
}

% ── Colors (match Module 1 lesson plan) ──────────────────────
\definecolor{stevens}{RGB}{163,31,52}
\definecolor{darkgray}{RGB}{60,60,60}
\definecolor{lightgray}{RGB}{240,240,240}
\definecolor{accentblue}{RGB}{30,80,140}

\hypersetup{
	colorlinks=true,
	linkcolor=stevens,
	urlcolor=accentblue,
	citecolor=darkgray,
}

% ── Defaults for header variables (handout overrides these) ──
\providecommand{\labid}{?}
\providecommand{\labtitleshort}{Lab Handout}

% ── Running header / footer ──────────────────────────────────
\pagestyle{fancy}
\fancyhf{}
\fancyhead[L]{\small\textcolor{darkgray}{Module 1 --- Lab Handout \labid}}
\fancyhead[R]{\small\textcolor{darkgray}{\labtitleshort}}
\fancyfoot[C]{\small\thepage}
\renewcommand{\headrulewidth}{0.4pt}

% ── Section formatting ───────────────────────────────────────
\titleformat{\section}{\Large\bfseries\color{stevens}}{\thesection}{0.6em}{}
\titleformat{\subsection}{\large\bfseries\color{darkgray}}{\thesubsection}{0.6em}{}

% ── Box styles ───────────────────────────────────────────────
\tcbset{
	infobox/.style={
		enhanced, colback=blue!3, colframe=accentblue, fonttitle=\bfseries,
		title={#1}, coltitle=white, breakable, left=6pt, right=6pt, top=4pt, bottom=4pt
	},
	warning/.style={
		enhanced, colback=red!3, colframe=stevens, fonttitle=\bfseries,
		title={#1}, coltitle=white, breakable, left=6pt, right=6pt, top=4pt, bottom=4pt
	},
	deliverable/.style={
		enhanced, colback=gray!5, colframe=darkgray, fonttitle=\bfseries,
		title={#1}, coltitle=white, breakable, left=6pt, right=6pt, top=4pt, bottom=4pt
	},
}

% ── Figure helper: always draws a placeholder box (images suppressed) ─
% Usage: \labfigure{filename}{width}{caption}{label}
\newcommand{\labfigure}[4]{%
	\begin{figure}[h!]
		\centering
		\fbox{\begin{minipage}[c][2.2in][c]{#2}
			\centering\color{darkgray}\small
				\textbf{[Figure placeholder]}\\[4pt]
				\texttt{\detokenize{#1}}\\[4pt]
				\textit{#3}
		\end{minipage}}%
		\caption{#3}
		\label{#4}
	\end{figure}%
}

% ── Title banner (call at top of document body) ──────────────
\newcommand{\labheader}[2]{%
	\begin{center}
		{\LARGE\bfseries\color{stevens} Lab Handout #1}\\[3pt]
		{\Large\bfseries #2}\\[6pt]
		{\small\color{darkgray} Capstone --- Module 1: \RF{} Hardware Foundations}
	\end{center}
	\vspace{4pt}
	{\color{darkgray}\rule{\textwidth}{0.4pt}}
	\vspace{10pt}
}

% BEFORE \begin{document}. Compiles with a single pdflatex pass (no glossaries tool).

\usepackage[T1]{fontenc}
\usepackage[margin=1in]{geometry}
\usepackage{titlesec}
\usepackage{enumitem}
\usepackage{booktabs}
\usepackage{tabularx}
\usepackage{graphicx}
\graphicspath{{Images/}{../Images/}}
\usepackage{xcolor}
\usepackage{hyperref}
\usepackage{fancyhdr}
\usepackage{amsmath}
\usepackage{amssymb}
\usepackage{tcolorbox}
\usepackage{caption}
\tcbuselibrary{skins,breakable}

% ── Glossary (per-document, built from shared glossary.tex) ──
% Uses \makenoidxglossaries so no external makeglossaries tool is
% needed: sorting and printing is done by LaTeX itself across two
% pdflatex passes (same cadence as the TOC). Only acronyms that
% are actually \gls{}-used in the handout appear in the printed
% glossary, and each \gls{} reference is a live hyperlink that
% jumps to its entry in the Acronyms section at the end.
\usepackage[acronym,nomain,nopostdot]{glossaries}
\makenoidxglossaries
% Shared acronym macro container for EE800/820 lesson plan modules.
% Usage: type \IDE, \FPGA, \UART, etc. First use is expanded; later uses show acronym only.
% Requires in each document preamble:
%   \usepackage[acronym,toc]{glossaries}
%   \makeglossaries
%   \input{../glossary}
% Print used acronyms near end of document:
%   \printglossary[type=\acronymtype,title={Acronyms Used}]

% Acronym definitions (only used entries appear in the printed glossary).
\newacronym{ask}{ASK}{Amplitude Shift Keying}
\newacronym{bga}{BGA}{Ball Grid Array}
\newacronym{bom}{BOM}{Bill of Materials}
\newacronym{bram}{BRAM}{Block Random Access Memory}
\newacronym{crc}{CRC}{Cyclic Redundancy Check}
\newacronym{css}{CSS}{Chirp Spread Spectrum}
\newacronym{dma}{DMA}{Direct Memory Access}
\newacronym{dsp}{DSP}{Digital Signal Processing}
\newacronym{eda}{EDA}{Exploratory Data Analysis}
\newacronym{fifo}{FIFO}{First In, First Out}
\newacronym{fpga}{FPGA}{Field-Programmable Gate Array}
\newacronym{fsk}{FSK}{Frequency Shift Keying}
\newacronym{fsm}{FSM}{Finite State Machine}
\newacronym{gpio}{GPIO}{General-Purpose Input/Output}
\newacronym{hal}{HAL}{Hardware Abstraction Layer}
\newacronym{i2c}{I\textsuperscript{2}C}{Inter-Integrated Circuit}
\newacronym{ide}{IDE}{Integrated Development Environment}
\newacronym{ila}{ILA}{Integrated Logic Analyzer}
\newacronym{ism}{ISM}{Industrial, Scientific, and Medical}
\newacronym{isr}{ISR}{Interrupt Service Routine}
\newacronym{lut}{LUT}{Look-Up Table}
\newacronym{mcu}{MCU}{Microcontroller Unit}
\newacronym{ml}{ML}{Machine Learning}
\newacronym{ook}{OOK}{On-Off Keying}
\newacronym{per}{PER}{Packet Error Rate}
\newacronym{rf}{RF}{Radio Frequency}
\newacronym{rssi}{RSSI}{Received Signal Strength Indicator}
\newacronym{rtl}{RTL}{Register-Transfer Level}
\newacronym{snr}{SNR}{Signal-to-Noise Ratio}
\newacronym{spi}{SPI}{Serial Peripheral Interface}
\newacronym{uart}{UART}{Universal Asynchronous Receiver/Transmitter}
\newacronym{vhdl}{VHDL}{VHSIC Hardware Description Language}

% Convenience wrappers so you can type \IDE, \RF, etc.
\providecommand{\ASK}{\gls{ask}}
\providecommand{\BGA}{\gls{bga}}
\providecommand{\BOM}{\gls{bom}}
\providecommand{\BRAM}{\gls{bram}}
\providecommand{\CRC}{\gls{crc}}
\providecommand{\CSS}{\gls{css}}
\providecommand{\DMA}{\gls{dma}}
\providecommand{\DSP}{\gls{dsp}}
\providecommand{\EDA}{\gls{eda}}
\providecommand{\FIFO}{\gls{fifo}}
\providecommand{\FPGA}{\gls{fpga}}
\providecommand{\FSK}{\gls{fsk}}
\providecommand{\FSM}{\gls{fsm}}
\providecommand{\GPIO}{\gls{gpio}}
\providecommand{\HAL}{\gls{hal}}
\providecommand{\IIC}{\gls{i2c}}
\providecommand{\IDE}{\gls{ide}}
\providecommand{\ILA}{\gls{ila}}
\providecommand{\ISM}{\gls{ism}}
\providecommand{\ISR}{\gls{isr}}
\providecommand{\LUT}{\gls{lut}}
\providecommand{\MCU}{\gls{mcu}}
\providecommand{\ML}{\gls{ml}}
\providecommand{\OOK}{\gls{ook}}
\providecommand{\PER}{\gls{per}}
\providecommand{\RF}{\gls{rf}}
\providecommand{\RSSI}{\gls{rssi}}
\providecommand{\RTL}{\gls{rtl}}
\providecommand{\SNR}{\gls{snr}}
\providecommand{\SPI}{\gls{spi}}
\providecommand{\UART}{\gls{uart}}
\providecommand{\VHDL}{\gls{vhdl}}


% Call \printhandoutglossary just before \end{document} in each handout.
\newcommand{\printhandoutglossary}{%
	\clearpage
	\printnoidxglossary[type=\acronymtype,title={Glossary}]%
}

% ── Colors (match Module 1 lesson plan) ──────────────────────
\definecolor{stevens}{RGB}{163,31,52}
\definecolor{darkgray}{RGB}{60,60,60}
\definecolor{lightgray}{RGB}{240,240,240}
\definecolor{accentblue}{RGB}{30,80,140}

\hypersetup{
	colorlinks=true,
	linkcolor=stevens,
	urlcolor=accentblue,
	citecolor=darkgray,
}

% ── Defaults for header variables (handout overrides these) ──
\providecommand{\labid}{?}
\providecommand{\labtitleshort}{Lab Handout}

% ── Running header / footer ──────────────────────────────────
\pagestyle{fancy}
\fancyhf{}
\fancyhead[L]{\small\textcolor{darkgray}{Module 1 --- Lab Handout \labid}}
\fancyhead[R]{\small\textcolor{darkgray}{\labtitleshort}}
\fancyfoot[C]{\small\thepage}
\renewcommand{\headrulewidth}{0.4pt}

% ── Section formatting ───────────────────────────────────────
\titleformat{\section}{\Large\bfseries\color{stevens}}{\thesection}{0.6em}{}
\titleformat{\subsection}{\large\bfseries\color{darkgray}}{\thesubsection}{0.6em}{}

% ── Box styles ───────────────────────────────────────────────
\tcbset{
	infobox/.style={
		enhanced, colback=blue!3, colframe=accentblue, fonttitle=\bfseries,
		title={#1}, coltitle=white, breakable, left=6pt, right=6pt, top=4pt, bottom=4pt
	},
	warning/.style={
		enhanced, colback=red!3, colframe=stevens, fonttitle=\bfseries,
		title={#1}, coltitle=white, breakable, left=6pt, right=6pt, top=4pt, bottom=4pt
	},
	deliverable/.style={
		enhanced, colback=gray!5, colframe=darkgray, fonttitle=\bfseries,
		title={#1}, coltitle=white, breakable, left=6pt, right=6pt, top=4pt, bottom=4pt
	},
}

% ── Figure helper: always draws a placeholder box (images suppressed) ─
% Usage: \labfigure{filename}{width}{caption}{label}
\newcommand{\labfigure}[4]{%
	\begin{figure}[h!]
		\centering
		\fbox{\begin{minipage}[c][2.2in][c]{#2}
			\centering\color{darkgray}\small
				\textbf{[Figure placeholder]}\\[4pt]
				\texttt{\detokenize{#1}}\\[4pt]
				\textit{#3}
		\end{minipage}}%
		\caption{#3}
		\label{#4}
	\end{figure}%
}

% ── Title banner (call at top of document body) ──────────────
\newcommand{\labheader}[2]{%
	\begin{center}
		{\LARGE\bfseries\color{stevens} Lab Handout #1}\\[3pt]
		{\Large\bfseries #2}\\[6pt]
		{\small\color{darkgray} Capstone --- Module 1: \RF{} Hardware Foundations}
	\end{center}
	\vspace{4pt}
	{\color{darkgray}\rule{\textwidth}{0.4pt}}
	\vspace{10pt}
}

% BEFORE \begin{document}. Compiles with a single pdflatex pass (no glossaries tool).

\usepackage[T1]{fontenc}
\usepackage[margin=1in]{geometry}
\usepackage{titlesec}
\usepackage{enumitem}
\usepackage{booktabs}
\usepackage{tabularx}
\usepackage{graphicx}
\graphicspath{{Images/}{../Images/}}
\usepackage{xcolor}
\usepackage{hyperref}
\usepackage{fancyhdr}
\usepackage{amsmath}
\usepackage{amssymb}
\usepackage{tcolorbox}
\usepackage{caption}
\tcbuselibrary{skins,breakable}

% ── Glossary (per-document, built from shared glossary.tex) ──
% Uses \makenoidxglossaries so no external makeglossaries tool is
% needed: sorting and printing is done by LaTeX itself across two
% pdflatex passes (same cadence as the TOC). Only acronyms that
% are actually \gls{}-used in the handout appear in the printed
% glossary, and each \gls{} reference is a live hyperlink that
% jumps to its entry in the Acronyms section at the end.
\usepackage[acronym,nomain,nopostdot]{glossaries}
\makenoidxglossaries
% Shared acronym macro container for EE800/820 lesson plan modules.
% Usage: type \IDE, \FPGA, \UART, etc. First use is expanded; later uses show acronym only.
% Requires in each document preamble:
%   \usepackage[acronym,toc]{glossaries}
%   \makeglossaries
%   % Shared acronym macro container for EE800/820 lesson plan modules.
% Usage: type \IDE, \FPGA, \UART, etc. First use is expanded; later uses show acronym only.
% Requires in each document preamble:
%   \usepackage[acronym,toc]{glossaries}
%   \makeglossaries
%   \input{../glossary}
% Print used acronyms near end of document:
%   \printglossary[type=\acronymtype,title={Acronyms Used}]

% Acronym definitions (only used entries appear in the printed glossary).
\newacronym{ask}{ASK}{Amplitude Shift Keying}
\newacronym{bga}{BGA}{Ball Grid Array}
\newacronym{bom}{BOM}{Bill of Materials}
\newacronym{bram}{BRAM}{Block Random Access Memory}
\newacronym{crc}{CRC}{Cyclic Redundancy Check}
\newacronym{css}{CSS}{Chirp Spread Spectrum}
\newacronym{dma}{DMA}{Direct Memory Access}
\newacronym{dsp}{DSP}{Digital Signal Processing}
\newacronym{eda}{EDA}{Exploratory Data Analysis}
\newacronym{fifo}{FIFO}{First In, First Out}
\newacronym{fpga}{FPGA}{Field-Programmable Gate Array}
\newacronym{fsk}{FSK}{Frequency Shift Keying}
\newacronym{fsm}{FSM}{Finite State Machine}
\newacronym{gpio}{GPIO}{General-Purpose Input/Output}
\newacronym{hal}{HAL}{Hardware Abstraction Layer}
\newacronym{i2c}{I\textsuperscript{2}C}{Inter-Integrated Circuit}
\newacronym{ide}{IDE}{Integrated Development Environment}
\newacronym{ila}{ILA}{Integrated Logic Analyzer}
\newacronym{ism}{ISM}{Industrial, Scientific, and Medical}
\newacronym{isr}{ISR}{Interrupt Service Routine}
\newacronym{lut}{LUT}{Look-Up Table}
\newacronym{mcu}{MCU}{Microcontroller Unit}
\newacronym{ml}{ML}{Machine Learning}
\newacronym{ook}{OOK}{On-Off Keying}
\newacronym{per}{PER}{Packet Error Rate}
\newacronym{rf}{RF}{Radio Frequency}
\newacronym{rssi}{RSSI}{Received Signal Strength Indicator}
\newacronym{rtl}{RTL}{Register-Transfer Level}
\newacronym{snr}{SNR}{Signal-to-Noise Ratio}
\newacronym{spi}{SPI}{Serial Peripheral Interface}
\newacronym{uart}{UART}{Universal Asynchronous Receiver/Transmitter}
\newacronym{vhdl}{VHDL}{VHSIC Hardware Description Language}

% Convenience wrappers so you can type \IDE, \RF, etc.
\providecommand{\ASK}{\gls{ask}}
\providecommand{\BGA}{\gls{bga}}
\providecommand{\BOM}{\gls{bom}}
\providecommand{\BRAM}{\gls{bram}}
\providecommand{\CRC}{\gls{crc}}
\providecommand{\CSS}{\gls{css}}
\providecommand{\DMA}{\gls{dma}}
\providecommand{\DSP}{\gls{dsp}}
\providecommand{\EDA}{\gls{eda}}
\providecommand{\FIFO}{\gls{fifo}}
\providecommand{\FPGA}{\gls{fpga}}
\providecommand{\FSK}{\gls{fsk}}
\providecommand{\FSM}{\gls{fsm}}
\providecommand{\GPIO}{\gls{gpio}}
\providecommand{\HAL}{\gls{hal}}
\providecommand{\IIC}{\gls{i2c}}
\providecommand{\IDE}{\gls{ide}}
\providecommand{\ILA}{\gls{ila}}
\providecommand{\ISM}{\gls{ism}}
\providecommand{\ISR}{\gls{isr}}
\providecommand{\LUT}{\gls{lut}}
\providecommand{\MCU}{\gls{mcu}}
\providecommand{\ML}{\gls{ml}}
\providecommand{\OOK}{\gls{ook}}
\providecommand{\PER}{\gls{per}}
\providecommand{\RF}{\gls{rf}}
\providecommand{\RSSI}{\gls{rssi}}
\providecommand{\RTL}{\gls{rtl}}
\providecommand{\SNR}{\gls{snr}}
\providecommand{\SPI}{\gls{spi}}
\providecommand{\UART}{\gls{uart}}
\providecommand{\VHDL}{\gls{vhdl}}


% Print used acronyms near end of document:
%   \printglossary[type=\acronymtype,title={Acronyms Used}]

% Acronym definitions (only used entries appear in the printed glossary).
\newacronym{ask}{ASK}{Amplitude Shift Keying}
\newacronym{bga}{BGA}{Ball Grid Array}
\newacronym{bom}{BOM}{Bill of Materials}
\newacronym{bram}{BRAM}{Block Random Access Memory}
\newacronym{crc}{CRC}{Cyclic Redundancy Check}
\newacronym{css}{CSS}{Chirp Spread Spectrum}
\newacronym{dma}{DMA}{Direct Memory Access}
\newacronym{dsp}{DSP}{Digital Signal Processing}
\newacronym{eda}{EDA}{Exploratory Data Analysis}
\newacronym{fifo}{FIFO}{First In, First Out}
\newacronym{fpga}{FPGA}{Field-Programmable Gate Array}
\newacronym{fsk}{FSK}{Frequency Shift Keying}
\newacronym{fsm}{FSM}{Finite State Machine}
\newacronym{gpio}{GPIO}{General-Purpose Input/Output}
\newacronym{hal}{HAL}{Hardware Abstraction Layer}
\newacronym{i2c}{I\textsuperscript{2}C}{Inter-Integrated Circuit}
\newacronym{ide}{IDE}{Integrated Development Environment}
\newacronym{ila}{ILA}{Integrated Logic Analyzer}
\newacronym{ism}{ISM}{Industrial, Scientific, and Medical}
\newacronym{isr}{ISR}{Interrupt Service Routine}
\newacronym{lut}{LUT}{Look-Up Table}
\newacronym{mcu}{MCU}{Microcontroller Unit}
\newacronym{ml}{ML}{Machine Learning}
\newacronym{ook}{OOK}{On-Off Keying}
\newacronym{per}{PER}{Packet Error Rate}
\newacronym{rf}{RF}{Radio Frequency}
\newacronym{rssi}{RSSI}{Received Signal Strength Indicator}
\newacronym{rtl}{RTL}{Register-Transfer Level}
\newacronym{snr}{SNR}{Signal-to-Noise Ratio}
\newacronym{spi}{SPI}{Serial Peripheral Interface}
\newacronym{uart}{UART}{Universal Asynchronous Receiver/Transmitter}
\newacronym{vhdl}{VHDL}{VHSIC Hardware Description Language}

% Convenience wrappers so you can type \IDE, \RF, etc.
\providecommand{\ASK}{\gls{ask}}
\providecommand{\BGA}{\gls{bga}}
\providecommand{\BOM}{\gls{bom}}
\providecommand{\BRAM}{\gls{bram}}
\providecommand{\CRC}{\gls{crc}}
\providecommand{\CSS}{\gls{css}}
\providecommand{\DMA}{\gls{dma}}
\providecommand{\DSP}{\gls{dsp}}
\providecommand{\EDA}{\gls{eda}}
\providecommand{\FIFO}{\gls{fifo}}
\providecommand{\FPGA}{\gls{fpga}}
\providecommand{\FSK}{\gls{fsk}}
\providecommand{\FSM}{\gls{fsm}}
\providecommand{\GPIO}{\gls{gpio}}
\providecommand{\HAL}{\gls{hal}}
\providecommand{\IIC}{\gls{i2c}}
\providecommand{\IDE}{\gls{ide}}
\providecommand{\ILA}{\gls{ila}}
\providecommand{\ISM}{\gls{ism}}
\providecommand{\ISR}{\gls{isr}}
\providecommand{\LUT}{\gls{lut}}
\providecommand{\MCU}{\gls{mcu}}
\providecommand{\ML}{\gls{ml}}
\providecommand{\OOK}{\gls{ook}}
\providecommand{\PER}{\gls{per}}
\providecommand{\RF}{\gls{rf}}
\providecommand{\RSSI}{\gls{rssi}}
\providecommand{\RTL}{\gls{rtl}}
\providecommand{\SNR}{\gls{snr}}
\providecommand{\SPI}{\gls{spi}}
\providecommand{\UART}{\gls{uart}}
\providecommand{\VHDL}{\gls{vhdl}}


% Call \printhandoutglossary just before \end{document} in each handout.
\newcommand{\printhandoutglossary}{%
	\clearpage
	\printnoidxglossary[type=\acronymtype,title={Glossary}]%
}

% ── Colors (match Module 1 lesson plan) ──────────────────────
\definecolor{stevens}{RGB}{163,31,52}
\definecolor{darkgray}{RGB}{60,60,60}
\definecolor{lightgray}{RGB}{240,240,240}
\definecolor{accentblue}{RGB}{30,80,140}

\hypersetup{
	colorlinks=true,
	linkcolor=stevens,
	urlcolor=accentblue,
	citecolor=darkgray,
}

% ── Defaults for header variables (handout overrides these) ──
\providecommand{\labid}{?}
\providecommand{\labtitleshort}{Lab Handout}

% ── Running header / footer ──────────────────────────────────
\pagestyle{fancy}
\fancyhf{}
\fancyhead[L]{\small\textcolor{darkgray}{Module 1 --- Lab Handout \labid}}
\fancyhead[R]{\small\textcolor{darkgray}{\labtitleshort}}
\fancyfoot[C]{\small\thepage}
\renewcommand{\headrulewidth}{0.4pt}

% ── Section formatting ───────────────────────────────────────
\titleformat{\section}{\Large\bfseries\color{stevens}}{\thesection}{0.6em}{}
\titleformat{\subsection}{\large\bfseries\color{darkgray}}{\thesubsection}{0.6em}{}

% ── Box styles ───────────────────────────────────────────────
\tcbset{
	infobox/.style={
		enhanced, colback=blue!3, colframe=accentblue, fonttitle=\bfseries,
		title={#1}, coltitle=white, breakable, left=6pt, right=6pt, top=4pt, bottom=4pt
	},
	warning/.style={
		enhanced, colback=red!3, colframe=stevens, fonttitle=\bfseries,
		title={#1}, coltitle=white, breakable, left=6pt, right=6pt, top=4pt, bottom=4pt
	},
	deliverable/.style={
		enhanced, colback=gray!5, colframe=darkgray, fonttitle=\bfseries,
		title={#1}, coltitle=white, breakable, left=6pt, right=6pt, top=4pt, bottom=4pt
	},
}

% ── Figure helper: always draws a placeholder box (images suppressed) ─
% Usage: \labfigure{filename}{width}{caption}{label}
\newcommand{\labfigure}[4]{%
	\begin{figure}[h!]
		\centering
		\fbox{\begin{minipage}[c][2.2in][c]{#2}
			\centering\color{darkgray}\small
				\textbf{[Figure placeholder]}\\[4pt]
				\texttt{\detokenize{#1}}\\[4pt]
				\textit{#3}
		\end{minipage}}%
		\caption{#3}
		\label{#4}
	\end{figure}%
}

% ── Title banner (call at top of document body) ──────────────
\newcommand{\labheader}[2]{%
	\begin{center}
		{\LARGE\bfseries\color{stevens} Lab Handout #1}\\[3pt]
		{\Large\bfseries #2}\\[6pt]
		{\small\color{darkgray} Capstone --- Module 1: \RF{} Hardware Foundations}
	\end{center}
	\vspace{4pt}
	{\color{darkgray}\rule{\textwidth}{0.4pt}}
	\vspace{10pt}
}


\begin{document}
\labheader{1-E}{\FPGA{} Toolchain Validation: Nexys A7 Loopback Design}

\section{Objective}
Validate the AMD Vivado toolchain and the Digilent Nexys A7-100T board by synthesizing, implementing, and programming a trivial switch-to-\mbox{LED} loopback design. This is the \FPGA{}-side counterpart to Lab Handouts~1-B/1-C and confirms that you can move on to \RTL{} development in Module~2.

\section{Prerequisites}
\begin{itemize}[leftmargin=2em]
	\item No prior lab handouts required for this one; it runs in parallel to the STM32 bring-up track.
	\item Basic \VHDL{} syntax familiarity (see Module~1 Prerequisites for references).
\end{itemize}

\section{Materials}
\begin{itemize}[leftmargin=2em]
	\item Digilent Nexys A7-100T development board (XC7A100T-1CSG324C).
	\item USB Micro-B cable.
	\item Host PC with $\geq$50\,GB free disk space for Vivado.
\end{itemize}

\section{Procedure}

\subsection{Step 1 --- Install Vivado (first-time only)}
\begin{enumerate}[leftmargin=2em]
	\item Download Vivado ML Standard (the free edition that covers Artix-7) from \href{https://www.xilinx.com/support/download.html}{xilinx.com/support/download.html}. Create a free AMD account if you don't already have one.
	\item Run the installer. When prompted for device families, select \textbf{Artix-7} and \textbf{7 Series} at minimum; deselecting everything else saves substantial disk space and install time.
	\item On Windows, let the installer install the cable drivers when prompted at the end.
	\item Launch Vivado once to verify the license. The free Standard edition does not require a license file for Artix-7 devices.
\end{enumerate}

\subsection{Step 2 --- Obtain the master constraint file}
\begin{enumerate}[leftmargin=2em]
	\item Download \texttt{Nexys-A7-100T-Master.xdc} from Digilent's Nexys A7 reference page (search: ``Digilent Nexys A7 master XDC'').
	\item Save it locally. You will add it to the project and uncomment only the pins you need.
\end{enumerate}

\subsection{Step 3 --- Create the Vivado project}
\begin{enumerate}[leftmargin=2em]
	\item Vivado $\to$ \textit{File} $\to$ \textit{Project} $\to$ \textit{New}. Name the project \texttt{nexys\_loopback}.
	\item Project type: \textbf{\mbox{RTL} Project}. Do not specify sources yet.
	\item Target device: \textbf{xc7a100tcsg324-1}.
	\item Finish the wizard.
\end{enumerate}

\subsection{Step 4 --- Add the \VHDL{} top module}
Create a new source file \texttt{loopback\_top.vhd} with the following contents:
\begin{verbatim}
library IEEE;
use IEEE.STD_LOGIC_1164.ALL;

entity loopback_top is
    Port ( sw  : in  STD_LOGIC_VECTOR(3 downto 0);
           led : out STD_LOGIC_VECTOR(3 downto 0));
end loopback_top;

architecture Behavioral of loopback_top is
begin
    led <= sw;
end Behavioral;
\end{verbatim}

\subsection{Step 5 --- Add and prune the constraint file}
\begin{enumerate}[leftmargin=2em]
	\item Add \texttt{Nexys-A7-100T-Master.xdc} as a constraints source.
	\item Open the file. Uncomment \textit{only} the lines for \texttt{sw[0]}--\texttt{sw[3]} (slide switches on pins J15, L16, M13, R15 --- double-check against the master file you downloaded) and \texttt{led[0]}--\texttt{led[3]} (pins H17, K15, J13, N14).
	\item Save.
\end{enumerate}

\subsection{Step 6 --- Synthesize, implement, generate bitstream}
\begin{enumerate}[leftmargin=2em]
	\item In the Flow Navigator, click \textbf{Run Synthesis}. Wait for it to complete with zero errors.
	\item Click \textbf{Run Implementation}. Wait for completion.
	\item Click \textbf{Generate Bitstream}. This produces \texttt{loopback\_top.bit}.
	\item Open the utilization report. You should see \mbox{LUT} usage in the single-digits out of 63{,}400 --- essentially nothing, which confirms the synthesizer optimized the direct assignment.
\end{enumerate}

\subsection{Step 7 --- Program the Nexys A7}
\begin{enumerate}[leftmargin=2em]
	\item Connect the Nexys A7 to the host PC via USB. Power the board on. The ``FPGA DONE'' \mbox{LED} will be off initially.
	\item Open the Hardware Manager in Vivado. Click \textit{Open Target} $\to$ \textit{Auto Connect}. Vivado should detect the XC7A100T.
	\item Right-click the device and select \textit{Program Device}. Point at the \texttt{loopback\_top.bit} file and click \textit{Program}.
	\item After a successful program, the ``FPGA DONE'' \mbox{LED} should illuminate.
\end{enumerate}

\subsection{Step 8 --- Verify physical behavior}
\begin{enumerate}[leftmargin=2em]
	\item Flip slide switch \texttt{SW0}. \mbox{LED} \texttt{LD0} should turn on.
	\item Repeat for \texttt{SW1}--\texttt{SW3} and \texttt{LD1}--\texttt{LD3}. Each switch must drive exactly one \mbox{LED}.
	\item Power-cycle the board. The bitstream is volatile --- the \mbox{LED}s should go dark and the DONE \mbox{LED} should turn off.
\end{enumerate}

\section{Verification Checklist}
\begin{itemize}[leftmargin=2em]
	\item[$\square$] Vivado project created targeting \texttt{xc7a100tcsg324-1}.
	\item[$\square$] Synthesis, implementation, and bitstream generation complete with zero errors.
	\item[$\square$] Bitstream programs successfully over the Nexys A7 USB-JTAG link.
	\item[$\square$] All four slide switches independently control their corresponding \mbox{LED}s.
	\item[$\square$] Power cycle clears the bitstream (expected for volatile configuration).
\end{itemize}

\begin{tcolorbox}[warning={Common Pitfalls}]
	\begin{itemize}[leftmargin=1.2em]
		\item \textbf{Constraint not found errors}: you edited the wrong \texttt{.xdc} file or added it as a generic source rather than a constraints source.
		\item \textbf{Hardware Manager does not see the board}: on Linux, \texttt{cable drivers} must be installed manually (see \texttt{<Vivado>/data/xicom/cable\_drivers/}).
		\item \textbf{Device programs but \mbox{LED}s don't respond}: the power switch on the Nexys A7 is off, or the board is powered from USB without the jumper set to USB --- check the \texttt{JP1} / \texttt{JP3} jumpers.
	\end{itemize}
\end{tcolorbox}

\begin{tcolorbox}[deliverable={Lab Handout 1-E Deliverable}]
	\begin{enumerate}[leftmargin=1.2em]
		\item \VHDL{} source file \texttt{loopback\_top.vhd}.
		\item Pruned \texttt{.xdc} constraints (just the 8 active lines).
		\item Screenshot of the Vivado utilization report after successful bitstream generation.
		\item Short video or photo sequence showing each switch independently controlling its \mbox{LED}.
	\end{enumerate}
\end{tcolorbox}

\end{document}
